\begin{frame} \ftx{Beispiel: Magnetische Flussdichte}
	\begin{bsp}{Magnetische Flussdichte}{}
		Im Inneren einer dicht gewickelten Ringspule soll die magnetische Feldstärke $H=100\,\frac{\mathrm{A}}{\mathrm{m}}$ erzeugt werden. Die Spule hat einen mittleren Radius von $5\,\mathrm{cm}$.
		\begin{enumerate}
			\item Berechnen Sie die erforderliche Stromstärke $I$ wenn die Spule mit $N=200$ Wicklungen versehen ist.\pause
			      
			      \s{Aus Gleichung \ref{Durchflutungsgesetz} und \ref{GlmagnFeldstaerke}:}
			      \begin{align*}
				      \varTheta & = H \cdot \ell_{\mathrm{m}} = N\cdot I                                                                                                                  \\
				      I         & =\frac{H\cdot\ell_{\mathrm{m}}}{N}=\frac{100\,\frac{\mathrm{A}}{\mathrm{m}}\cdot 2\cdot \pi\cdot 5\cdot 10^{-2}\,\mathrm{m}}{200} = 157,08\,\mathrm{mA}
			      \end{align*}\pause
			\item Wie groß wird die Flussdichte $B$ im Falle einer Luftspule ($\mu_{\mathrm{r}}=1$) oder einer eisengefüllten Spule ($\mu_{\mathrm{r}}=2000$ im Arbeitspunkt)?\pause
			      
			      \s{Aus Gleichung \ref{GlFlussdichte}:}
			      \begin{align*}
				      B_\mathrm{Luft}  & = \mu_0\cdot \mu_{\mathrm{r}}\cdot H = 1,256\cdot10^{-6}\,\tfrac{\mathrm{Vs}}{\mathrm{Am}} \cdot 100\,\tfrac{\mathrm{A}}{\mathrm{m}} = 125,6\,\mu\mathrm{T} \\
				      B_\mathrm{Eisen} & = 2000\cdot B_\mathrm{Luft} = 251,2\,\mathrm{mT}
			      \end{align*}
		\end{enumerate}
	\end{bsp}
\end{frame}

\section{Der magnetische Widerstand}

\begin{frame} \ftx{Magnetischer Widerstand}\index{Magnetischer Widerstand}

	\s{
		Im Kapitel \ref{durchflutung} wurde bereits die magnetische Spannung thematisiert, die auch Durchflutung genannt wird, sowie im Kapitel \ref{Flussdichte} der magnetische Fluss, der das Äquivalent zum elektrischen Strom darstellt. Es liegt daher die Vermutung nahe, dass es in einem magnetischen Kreis äquivalent zum ohmschen Widerstand auch einen magnetischen Widerstand\index{Magnetischer Widerstand} gibt. Er hat das Formelzeichen $R_{\mathrm{m}}$ und die Einheit $\frac{\mathrm{A}}{\mathrm{V}\mathrm{s}}$ und wird auch Reluktanz \index{Reluktanz}genannt.

		Ein Beispiel für einen magnetischen Kreis ist in Abbildung \ref{AbbMagnKreis} durch einen einfachen Eisenring mit einer einseitigen Spule dargestellt. Die Spule erzeugt eine Durchflutung $\varTheta$ analog zur elektrischen Spannung. Der Eisenkreis besteht aus vier Teilwiderständen, da der magnetische Widerstand sowohl von der Länge als auch von dem durchflossenen Querschnitt abhängig ist. Die vier Widerstände werden vom magnetischen Fluss $\varPhi$ durchströmt.
	}

	\fu{
		\input{Tikz/src/magnetischer_kreis}\alttext{Links im Bild ist ein magnetischer Kreis abgebildet, der aus einem Eisenring besteht. Dieser Ring ist von den vier magnetischen Teilwiderständen durchbrochen, die die vier Seiten des Eisenrings als darstellen. Diese Teilwiderstände sind abhängig von der Länge, dem Material und dem Querschnitt des Eisenrings. Rechts im Bild befindet sich die Abbildung eines Ersatzschaltbildes zu diesem Aufbau. Es zeigt, wie die magnetischen Teilwiderstände sowie der magnetische Fluss miteinander in Verbindung stehen, wenn eine Spannung an die Spule um den Eisenring angelegt wird und damit der magnetische Fluss die einzelnen Widerstände der vier Seiten des Eisenkerns nacheinander durchströmt.}
	}{{\bf Magnetischer Kreis mit einem Eisenring.} Die vier magnetischen Teilwiderstände sind abhängig von der Länge, vom Material und dem Querschnitt des durchflossenen Eisenrings. \label{AbbMagnKreis}}

	\s{In einfachen Anordnungen (wie beispielsweise in Abbildung \ref{AbbMagnKreis} zu sehen) kann unter Vernachlässigung der Ecken der Widerstand durch die Gleichung \ref{GlmagnWiderstand} ausgedrückt werden. $\ell_{\mathrm{m}}$ ist wie bei der magnetischen Feldstärke die mittlere Feldlinienlänge des Widerstandes innerhalb des magnetischen Kreises. Zur Ermittlung der mittleren Feldlininenlänge ${\ell_{\mathrm{m}}}$ werden die Längen aller Seiten addiert.
	$A$ ist die Querschnittsfläche, die vom magnetischen Fluss durchflossen wird. $\mu$ ist die Permeabilität\index{Permeabilität} des Materials.
	Der magnetische Widerstand $R_{\mathrm{m}}$ berechnet sich nun aus der aufsummierten Länge der mittleren Feldlinienlänge ${\ell_{\mathrm{m}}}$ geteilt durch das Produkt der durchflossenen Querschnittsfläche $A$ und der Permeabilität $\mu_{\mathrm{r}} \cdot \mu_0$.
	
	\begin{eq}
		R_{\mathrm{m}} =\frac{\ell_{\mathrm{m}}}{\mu_{\mathrm{r}}\cdot\mu_0\cdot A} \qquad\left[\frac{\mathrm{A}}{\mathrm{V}\cdot\mathrm{s}}\right]\label{GlmagnWiderstand} 
	\end{eq}
	
	\begin{Merksatz}{Magnetischer Widerstand}
		Für die Berechnung des magnetischen Widerstands $R_{\mathrm{m}}$ wird die mittlere Feldlinienlängen ${\ell_{\mathrm{m}}}$ durch das Produkt aus der Permeabilität des Materials ${\mu_{\mathrm{r}}\cdot\mu_0}$ und der Querschnittsfläche $A$ geteilt. Bei unterschiedlicher Materialbeschaffenheit oder unterschiedlicher Querschnittsflächen innerhalb des magnetisierten Körpers werden zunächst die Teilwiderstände errechnet und anschließenden zum Gesamtwiderstand aufsummiert.
	\end{Merksatz}
	}
	\pause	
	\b{\begin{minipage}{0.5\textwidth}}
	\b{
		\begin{eq}
			R_{\mathrm{m}} =\frac{\ell_{\mathrm{m}}}{\mu_{\mathrm{r}}\cdot\mu_0\cdot A} \qquad\left[\frac{\mathrm{A}}{\mathrm{V} \cdot \mathrm{s}}\right]
		\end{eq}}
	\b{\end{minipage}}%
\end{frame}

\subsection{Der magnetische Kreis}
\begin{frame} \ftx{Magnetischer Kreis}\index{Magnetischer Kreis}		
	\b{
		Analog zum elektrischen Ohmschen Gesetz gilt für den magnetischen Kreis:
		\begin{eqa}
			\varTheta &= R_{\mathrm{m}}\cdot \varPhi
		\end{eqa}
	}
	\s{Werden nun alle bekannten magnetischen Größen in ihrer Zusammenwirkung betrachtet, ergibt sich die magnetische Analogie zum Ohmschen Gesetz. Sie beschreibt, dass die magnetische Spannung $\varTheta$  dem Produkt aus dem magnetischen Widerstand $R_{\mathrm{m}}$ und dem magnetischen Fluss $\varPhi$ entspricht.	
	\begin{eqa}
		\varTheta &= R_{\mathrm{m}}\cdot \varPhi
	\end{eqa}
	\begin{Merksatz}{Zusammenhänge der magnetischen Feldgrößen}

		Die magnetische Feldstärke $H$ im Spezialfall einer Ringkernspule errechnet sich aus dem Produkt der Windungsanzahl $N$ und dem Stromfluss $I$, geteilt durch die mittlere Feldlinienlänge $\ell_{\mathrm{m}}$. \par
		\begin{equation*}
			H=\frac{N\cdot I}{\ell_{\mathrm{m}}}
		\end{equation*}		
		Die magnetische Flussdichte $B$ wird durch  Multiplikation der magnetischen Feldstärke $H$ mit der Permeabilität $\mu_{\mathrm{r}} \cdot \mu_0$ bestimmt. \par
		\begin{equation*}	
			B = \mu_{\mathrm{r}} \cdot \mu_0 \cdot H 
		\end{equation*}	

		Der magnetische Fluss $\varPhi$ ergibt sich aus dem Integral der magnetischen Flussdichte $B$ über einer Fläche $A$.		\begin{equation*} 
			\varPhi = \iint_A \vec{B} \cdot \mathrm{d} \vec{A}  
		\end{equation*}	

	\end{Merksatz}
	}
\end{frame}

